\section{Experiments}
\label{sec:exp}

We answer four questions.
Does \method improve closed-loop success over existing ultra-low-bit quantizers at the same bits per weight, across WAMs (\cref{sec:exp:main})?
Does the improvement hold on a real robot (\cref{sec:exp:real})?
What does each component contribute, and is the future term necessary (\cref{sec:exp:ablation})?
How much memory does it save, and how fast does it run (\cref{sec:exp:efficiency})?

\subsection{Experimental setup}
\label{sec:exp:setup}

\paragraph{Models.}
We quantize Fast-WAM~\cite{yuan2026fastwam} (6B parameters, of which 1B form the action expert) and \XX \TODO{second WAM, \eg Cosmos-Policy~\cite{kim2026cosmospolicy}}.
Fast-WAM first denoises the future latent in 10 steps and then denoises a 32-step action chunk in 10 steps while attending to it; the robot executes the first 10 actions and replans.
For a model that denoises the future and the action jointly, the reader starts from the future state at an intermediate denoising step and runs the remaining steps to the action.
Each model is compared with its own BF16 reference.

\paragraph{Benchmarks.}
We evaluate in simulation on LIBERO-Plus~\cite{fei2026liberoplus}, which perturbs the four LIBERO suites~\cite{liu2023libero} along seven dimensions.
We draw 840 episodes, 30 for each suite and perturbation dimension, and run every method on the same variants and seeds, so that episodes are paired across methods.
Real-world experiments are described in \cref{sec:exp:real}.

\paragraph{Baselines.}
We compare with scalar GPTQ~\cite{frantar2023gptq} and with the vector, lattice and trellis quantizers GPTVQ~\cite{vanbaalen2024gptvq}, QuIP\#~\cite{tseng2024quipsharp} and QTIP~\cite{tseng2024qtip}, at 2 and 3 bits per weight.
As is common, all methods keep the input and output interface layers (0.01\% of the weights) in BF16 and store the timestep-conditioning layers, whose inputs depend only on the denoising step, as exact per-step tables; both are counted in bits per weight.
All methods use the same calibration data and inference runtime.

\paragraph{Metrics.}
We report closed-loop success rates per suite and on average.
For efficiency, we report weight storage, peak memory and the latency of one policy call.

\paragraph{Calibration and implementation.}
All methods calibrate on the same 80 LIBERO contexts.
\method uses $\lambda = \XX$ and $\beta = \XX$ (\cref{sec:supp:method}) and 8 calibration updates.
Quantization runs on NVIDIA A100 GPUs; further details are in \cref{sec:supp:eval}.

\subsection{Main results}
\label{sec:exp:main}

\begin{table*}[t]
  \centering
  \small
  \setlength{\tabcolsep}{6pt}
  \begin{tabular}{@{}llccccccc@{}}
    \toprule
    Model & Method & bpw & Weights (GB) & Spatial & Object & Goal & Long & Avg. \\
    \midrule
    \multirow{11}{*}{Fast-WAM}
      & BF16    & 16   & 12.0 & 81.4 & 82.9 & 67.6 & 78.6 & 77.6 \\
    \cmidrule(l){2-9}
      & GPTQ    & \XX  & \XX  & \XX  & \XX  & \XX  & \XX  & \XX \\
      & GPTVQ   & \XX  & \XX  & \XX  & \XX  & \XX  & \XX  & \XX \\
      & QuIP\#  & \XX  & \XX  & \XX  & \XX  & \XX  & \XX  & \XX \\
      & QTIP    & 2.98 & \XX  & \XX  & \XX  & \XX  & \XX  & \XX \\
      & \cellcolor{gray!12}\method & \cellcolor{gray!12}2.98 & \cellcolor{gray!12}\XX & \cellcolor{gray!12}\XX & \cellcolor{gray!12}\XX & \cellcolor{gray!12}\XX & \cellcolor{gray!12}\XX & \cellcolor{gray!12}\XX \\
    \cmidrule(l){2-9}
      & GPTQ    & \XX  & \XX  & \XX  & \XX  & \XX  & \XX  & \XX \\
      & GPTVQ   & \XX  & \XX  & \XX  & \XX  & \XX  & \XX  & \XX \\
      & QuIP\#  & \XX  & \XX  & \XX  & \XX  & \XX  & \XX  & \XX \\
      & QTIP    & 1.99 & \XX  & \XX  & \XX  & \XX  & \XX  & \XX \\
      & \cellcolor{gray!12}\method & \cellcolor{gray!12}1.99 & \cellcolor{gray!12}\XX & \cellcolor{gray!12}\XX & \cellcolor{gray!12}\XX & \cellcolor{gray!12}\XX & \cellcolor{gray!12}\XX & \cellcolor{gray!12}\XX \\
    \midrule
    \multirow{11}{*}{\XX}
      & BF16    & 16   & \XX  & \XX  & \XX  & \XX  & \XX  & \XX \\
    \cmidrule(l){2-9}
      & GPTQ    & \XX  & \XX  & \XX  & \XX  & \XX  & \XX  & \XX \\
      & GPTVQ   & \XX  & \XX  & \XX  & \XX  & \XX  & \XX  & \XX \\
      & QuIP\#  & \XX  & \XX  & \XX  & \XX  & \XX  & \XX  & \XX \\
      & QTIP    & \XX  & \XX  & \XX  & \XX  & \XX  & \XX  & \XX \\
      & \cellcolor{gray!12}\method & \cellcolor{gray!12}\XX & \cellcolor{gray!12}\XX & \cellcolor{gray!12}\XX & \cellcolor{gray!12}\XX & \cellcolor{gray!12}\XX & \cellcolor{gray!12}\XX & \cellcolor{gray!12}\XX \\
    \cmidrule(l){2-9}
      & GPTQ    & \XX  & \XX  & \XX  & \XX  & \XX  & \XX  & \XX \\
      & GPTVQ   & \XX  & \XX  & \XX  & \XX  & \XX  & \XX  & \XX \\
      & QuIP\#  & \XX  & \XX  & \XX  & \XX  & \XX  & \XX  & \XX \\
      & QTIP    & \XX  & \XX  & \XX  & \XX  & \XX  & \XX  & \XX \\
      & \cellcolor{gray!12}\method & \cellcolor{gray!12}\XX & \cellcolor{gray!12}\XX & \cellcolor{gray!12}\XX & \cellcolor{gray!12}\XX & \cellcolor{gray!12}\XX & \cellcolor{gray!12}\XX & \cellcolor{gray!12}\XX \\
    \bottomrule
  \end{tabular}
  \caption{Closed-loop success rates (\%) on LIBERO-Plus over 840 paired episodes per method.
  Within each model, the upper group is at 3 bits per weight and the lower group at 2 bits.
  \textbf{Bold}: best quantized result in a group; \underline{underline}: second best.
  At 2 bits, \method is the only method within \XX points of BF16 on both models.}
  \label{tab:main}
\end{table*}

\cref{tab:main} reports closed-loop success on both models at 2 and 3 bits per weight.

\paragraph{Results on Fast-WAM.}
At 2 bits, \method reaches \XX\% average success, \XX points above QTIP, the strongest baseline, and recovers \XX\% of the gap between QTIP and BF16.
Scalar GPTQ collapses at this bit width, while the vector, lattice and trellis quantizers stay between \XX\% and \XX\%.
At 3 bits, all VQ baselines are within \XX points of BF16, and \method reaches \XX\%.
The gains concentrate on the \XX and \XX perturbations, where the baselines lose the most (\cref{sec:supp:results}).

\paragraph{Results on \XX.}
\XX.

\subsection{Real-world deployment}
\label{sec:exp:real}

\paragraph{Setup.}
We deploy Fast-WAM on an AgileX PiPER arm (6 degrees of freedom and a parallel gripper).
We fine-tune it on teleoperation data collected with LeRobot~\cite{cadene2024lerobot} from a third-person and a wrist camera, the same two-view layout as the LIBERO model, with about 80 demonstrations per task.
We quantize the fine-tuned BF16 model with calibration contexts drawn from the demonstrations and re-estimate the reader metric on it.
We compare BF16, 2-bit QTIP and 2-bit \method with \XX trials per task and method, using the same marked initial placements for all methods and randomizing the order of methods.

\paragraph{Tasks.}
\emph{Two-object pick-and-place}: put a can and a box into a basket, one after the other; this long-horizon task tests whether the quantized model completes a multi-step sequence.
\emph{Pouring}: pour beads from one cup into another, a continuous, dynamics-dependent motion.
\emph{Cloth folding}: fold a small cloth in half, which involves a deformable object.
Success criteria are listed in \cref{sec:supp:real}.

\begin{table}[t]
  \centering
  \small
  \setlength{\tabcolsep}{3pt}
  \begin{tabular}{@{}lcccccc@{}}
    \toprule
    & & \multicolumn{2}{c}{Pick-and-place} & & & \\
    \cmidrule(lr){3-4}
    Method & bpw & First & Both & Pouring & Folding & Avg. \\
    \midrule
    BF16    & 16   & \XX/\XX & \XX/\XX & \XX/\XX & \XX/\XX & \XX\% \\
    QTIP    & 1.99 & \XX/\XX & \XX/\XX & \XX/\XX & \XX/\XX & \XX\% \\
    \method & 1.99 & \XX/\XX & \XX/\XX & \XX/\XX & \XX/\XX & \XX\% \\
    \bottomrule
  \end{tabular}
  \caption{Real-world results on the AgileX PiPER arm (successes / trials).
  For pick-and-place, \emph{First} counts trials in which the first object reaches the basket and \emph{Both} counts full successes.
  The average uses full successes of the three tasks.}
  \label{tab:real}
\end{table}

\begin{figure}[t]
  \centering
  \figbox{Placeholder for Fig.~4 (real-world setup and rollouts).\\[2pt]
  Left: PiPER arm with third-person and wrist cameras.\\
  Right: key frames of the three tasks; one pick-and-place trial where QTIP fails and \method succeeds.}{0.55\linewidth}
  \caption{Real-world setup and example rollouts of 2-bit \method on the AgileX PiPER arm.}
  \label{fig:real}
\end{figure}

\paragraph{Results.}
\cref{tab:real} and \cref{fig:real} summarize the real-world results.
2-bit \method succeeds in \XX of \XX trials, compared with \XX for QTIP and \XX for BF16.
On pick-and-place, \XX.
The quantized model stores \XX\,GB of weights instead of 12\,GB.

\subsection{Ablation study}
\label{sec:exp:ablation}

\begin{table}[t]
  \centering
  \small
  \setlength{\tabcolsep}{4pt}
  \begin{tabular}{@{}llc@{}}
    \toprule
    & Variant & Avg. (\%) \\
    \midrule
    \multirow{3}{*}{Components}
      & QTIP (uniform bits)                     & \XX \\
      & + future-aware allocation (ii)          & \XX \\
      & + reader metric and calibration (i, iv) & \XX \\
    \midrule
    \multirow{4}{*}{Objective}
      & $A$                                      & \XX \\
      & $A + \lambda Z$                          & \XX \\
      & $A + \lambda_{\bM} Z_{\bM}$              & \XX \\
      & $A + \lambda_{\bM} Z_{\bM} + \beta K$    & \XX \\
    \midrule
    \multirow{4}{*}{Allocation}
      & Uniform                                  & \XX \\
      & Reversed (bits to the reader)            & \XX \\
      & Random layers within groups              & \XX \\
      & By future impact                         & \XX \\
    \midrule
    \multirow{3}{*}{Metric}
      & $\bM = \mathbf{I}$                       & \XX \\
      & Diagonal                                 & \XX \\
      & Low-rank + diagonal + $\alpha\mathbf{I}$ & \XX \\
    \bottomrule
  \end{tabular}
  \caption{Ablations on Fast-WAM at 2 bits per weight (LIBERO-Plus, \XX paired episodes).
  The last row of each block is the full \method.
  All variants use the same bit budget.}
  \label{tab:ablation}
\end{table}

\cref{tab:ablation} isolates each component on Fast-WAM at 2 bits per weight with the bit budget held fixed.

\paragraph{Each component contributes.}
Future-aware allocation raises the success rate of QTIP from \XX\% to \XX\%, and the reader metric with shared-table calibration adds \XX points.

\paragraph{Future fidelity is necessary.}
Calibrating the tables on the action alone reaches \XX\%.
Adding the future term raises it to \XX\% with the uniform metric and to \XX\% with the reader metric, and the reader error $K$ adds \XX points.

\paragraph{Allocating by future impact matters.}
At the same budget, uniform bits reach \XX\% and reversing the allocation toward the reader reaches \XX\%, so the gain comes from where the bits go rather than from mixed precision itself.
Choosing layers at random within the groups reaches \XX\%.

\paragraph{Metric form.}
The full reader metric reaches \XX\%, compared with \XX\% for its diagonal alone and \XX\% for the identity.

\subsection{Efficiency}
\label{sec:exp:efficiency}

\begin{table}[t]
  \centering
  \small
  \setlength{\tabcolsep}{3pt}
  \begin{tabular}{@{}lccccc@{}}
    \toprule
    Method & bpw & Weights & Peak mem. & Latency & Quant. time \\
           &     & (GB)    & (GB)      & (ms)    & (GPU-h) \\
    \midrule
    BF16    & 16   & 12.0 & \XX & \XX & -- \\
    QTIP    & 1.99 & \XX  & \XX & \XX & \XX \\
    \method & 1.99 & \XX  & \XX & \XX & \XX \\
    \bottomrule
  \end{tabular}
  \caption{Memory and latency of one Fast-WAM policy call on \XX.
  Quantization time is the total over all stages.}
  \label{tab:efficiency}
\end{table}

\cref{tab:efficiency} reports the cost of 2-bit \method on Fast-WAM.
It reduces weight storage from 12\,GB to \XX\,GB and peak memory from \XX\,GB to \XX\,GB, with a policy-call latency of \XX\,ms compared with \XX\,ms for BF16.
Its storage and latency match QTIP at the same bits per weight, since the two share the decoder and format.
Quantization takes \XX GPU-hours, of which the reader metric and calibration take \XX.
